\chapter{Trascrizione non utilizzatori}\label{app:risultati-interviste-non-utilizzatori}
In questa appendice sono riportate le trascrizioni delle interviste effettuate 
ai non utilizzatori del sistema di vendite attuale.
Le interviste hanno avuto l'obiettivo di raccogliere informazioni sulle esigenze odierne
 e future del sistema di preventivazione..
\section{Ruoli coinvolti e utilizzo attuale degli strumenti}
I reparti tecnici, di gestione commessa e logistici interagiscono in modo indiretto
o marginale con l'attuale database Microsoft Access, basandosi prevalentemente su passaggi cartacei o file esterni:
\begin{itemize}
      \item \textbf{Ufficio Project Management e Project Engineering:} Coordinano le commesse tra commerciale
            e reparti tecnici, definendo le specifiche e supervisionando installazione e post-vendita.
            Non utilizzano Access; consultano i contratti in formato cartaceo o scansionato per revisionare le
            specifiche tecniche e contattare i clienti.
      \item \textbf{Ufficio Tecnico Meccanico:} Si occupa della progettazione meccanica e dei disegni costruttivi.
            Non utilizza Access; recepisce le specifiche dal cartaceo tramite i \ac{PE} per verificare revisioni,
            limiti dei componenti e condizioni del sito di montaggio.
      \item \textbf{Ufficio Tecnico Elettrico:} Gestisce la progettazione di quadri e impianti elettrici.
            Attualmente sopperisce alle carenze del sistema utilizzando fogli Excel esterni
            per il dimensionamento/calcolo costi
            (\textit{QUAEL.xltm}) e per la validazione tecnica dei preventivi (\textit{PROGETTO CEPI - Scheda progetto r2.xlsx}).
      \item \textbf{Ufficio Logistica:} Gestisce spedizioni, dogane e trasporti. Utilizza Access solo marginalmente
            per verificare
            le date di consegna o correggere a valle indirizzi, causali, dati di trasporto e Incoterms.
            L'operatività ruota attorno
            a un
            ``File Logistico'' in Excel centralizzato e alla gestione di copie cartacee nelle cartelle di rete.
      \item \textbf{Ufficio Manualistica e Ufficio Acquisti:} La Manualistica redige manuali e certificazioni dei prodotti venduti;
            l'Ufficio Acquisti gestisce gli ordini ai fornitori e controlla i contratti commerciali per verificare i ricambi in
            assistenza ed eventuali discrepanze tecniche o commerciali.
\end{itemize}

\section{Criticità trasversali rilevate}

\subsection{Frammentazione documentale e supporto cartaceo}
\begin{itemize}
      \item \textbf{Disallineamento delle informazioni:} L'utilizzo di copie cartacee multiple (distribuite tra Vendite,
            Ufficio Tecnico, Elettrico e Logistica) e di annotazioni manuali sul cartaceo comporta la perdita di informazioni e
            la circolazione di versioni obsolete.
      \item \textbf{Duplicazione per tempistiche differenti:} Per le commesse miste (meccaniche ed elettriche),
            la Logistica è costretta a duplicare fisicamente le copie tecniche delle \ac{RA} per gestire le diverse
            tempistiche dei due uffici.
\end{itemize}

\subsection{Disconnessione anagrafica ed errori di input}
\begin{itemize}
      \item \textbf{Disallineamento Access vs Ad Hoc:} Access presenta anagrafiche spesso non aggiornate e prive
            di storicizzazione, mentre l'ERP Ad Hoc contiene i dati fiscali corretti (Sede Legale) ma risulta
            impreciso sulle sedi operative o di spedizione.
      \item \textbf{Propagazione degli errori:} La prassi commerciale di duplicare vecchie commesse/RA trascina errori
            pregressi che devono essere corretti manualmente dalla Logistica a valle del processo.
      \item \textbf{Omissione dati logistici:} Nelle offerte mancano quasi sempre informazioni essenziali per il
            trasporto quali Incoterms, tipologia di imballo, pesi e volumi.
\end{itemize}

\subsection{Gestione post-vendita e disallineamento per le dogane}
\begin{itemize}
      \item \textbf{Variazioni non tracciate:} Il metodo attuale punta alla vendita rapida,
            ma non traccia in modo strutturato i componenti aggiunti o modificati in fase di progettazione e realizzazione.
      \item \textbf{Criticità sulle Packing List doganali:} I commerciali vendono ``macro-codici'' (es. impianto completo),
            ma le dogane (specialmente extra-UE) esigono il dettaglio di ogni singolo pallet. Poiché il sistema non
            esplode i macro-codici nei colli fisici, l'associazione tra voci contrattuali, distinta tecnica e pallet preparati
            dagli spedizionieri avviene manualmente, generando un alto rischio di errori doganali.
\end{itemize}

\section{Requisiti funzionali emersi}

\subsection{Struttura del contratto, distinta base e codifica}
\begin{itemize}
      \item \textbf{Gestione per macro-codici e standardizzazione:}
            \begin{itemize}
                  \item Utilizzo di ``macro-codici'' contenitore in fase di vendita (es. \textit{Macro-codice Quadro},
                        \textit{Pannello Operatore}, \textit{TVB generico}) direttamente agganciati alla distinta base
                        (completa al 90\%), che l'Ufficio Tecnico completerà poi con i disegni specifici.
                  \item Spinta verso la ``standardizzazione a catalogo'' tramite l'analisi statistica dei \textit{Best Sellers}
                        (es. TVB standard, stellari taglie 1000/1500/2000), facendo coincidere il codice di vendita con quello di realizzazione
                        e separando il flusso standard da quello ``fuori standard''.
                  \item Allineamento univoco dei codici articolo tra il nuovo configuratore e il gestionale Ad Hoc.
            \end{itemize}
      \item \textbf{Separazione ed esplosione delle voci (Meccanica vs Elettrica):}
            \begin{itemize}
                  \item Creazione di due sezioni contrattuali ben distinte per i componenti meccanici e quelli elettrici,
                        superando la confusione attuale.
                  \item Esplosione automatica: l'inserimento di utenze meccaniche (es. motori, vibratori) deve generare
                        automaticamente le voci elettriche correlate nel contratto tecnico (es. gli inverter devono essere integrati
                        nel quadro o associati all'utenza specifica, non trattati come voci isolate).
            \end{itemize}
      \item \textbf{Rintracciabilità e collegamento logistico:}
            \begin{itemize}
                  \item Superamento dei codici generici per garantire la tracciabilità dell'esatta variante del
                        componente venduto (es. variante specifica di un \textit{WP104}).
                  \item Collegamento nativo tra distinta commerciale e distinta tecnica (tramite Rif CC)
                        fin dalla progettazione, per automatizzare l'associazione con i pallet fisici e la generazione delle Packing List.
            \end{itemize}
\end{itemize}

\subsection{Configuratore tecnico, normativo e logistico}
Il sistema deve guidare il commerciale attraverso vincoli all'ingresso (menu a tendina) e controlli automatici
di completezza prima del salvataggio:
\begin{itemize}
      \item \textbf{Guida e vincoli per la parte elettrica:}
            \begin{itemize}
                  \item Suggerimento automatico del materiale del quadro in base alle utenze inserite (es. taglia corretta
                        dell'inverter in base ai motori venduti) e domande finali per proporre accessori consigliati.
                  \item Inserimento automatico di default dei sezionatori a bordo macchina per ogni motore.
                  \item Distinzione netta tra ``quadro nuovo'' e ``ampliamento'': in caso di ampliamento, il sistema deve
                        indicare la disponibilità di spazio fisico nel quadro esistente per evitare errori di valutazione e costi imprevisti.
                  \item Campi tecnici obbligatori: Paese di destinazione, alimentazione (es. 230V 50Hz, 400V 60Hz, DC/AC),
                        materiale del quadro (inox, verniciato), tipologia canaline (chiusa, forata), numero/tipologia utenze e
                        certificazioni (es. UL/USA, la cui selezione deve aggiornare in automatico componenti e costi).
            \end{itemize}
      \item \textbf{Specifiche per manualistica, meccanica e acquisti:}
            \begin{itemize}
                  \item Campi dedicati alle certificazioni richieste per i componenti venduti (CE, ATEX, FDA, ecc.)
                        e alle norme di riferimento per i prodotti acquistati da terzi (ATEX, alimentari, ecc.).
                  \item Valorizzazione dettagliata delle caratteristiche tecniche (dati elettrici dei motori, dati
                        meccanici, allestimenti, accessori) per agevolare la stesura dei manuali.
                  \item Verifica delle condizioni del luogo di montaggio per attivare controlli specifici sulle macchine
                        (es. Globosilo esterni in zone sismiche), definire i trasporti e rilevare limitazioni d'uso dei componenti.
            \end{itemize}
      \item \textbf{Vincoli e alert logistici/doganali per Paese:}
            \begin{itemize}
                  \item Compilazione obbligatoria e guidata di Incoterms (DAP, DDP, FCA, ecc.), tipologia di imballo, pesi e volumi.
                  \item Generazione di alert automatici all'inserimento dell'ordine in base al Paese di destinazione per requisiti
                        di fumigazione/ISPM 15 sui pallet in legno (es. Brasile, Australia, Turchia), materiali proibiti
                        (es. ancoranti chimici vietati) e necessità di permessi eccezionali o montaggi.
            \end{itemize}
\end{itemize}

\subsection{Preventivazione, calcolo costi e varianti}
\begin{itemize}
      \item \textbf{Internalizzazione dei fogli di calcolo:} Integrazione nel software delle logiche di dimensionamento,
            calcolo costi e validazione preventiva oggi gestite sui file Excel esterni (\textit{QUAEL.xltm}
            e \textit{Scheda progetto r2.xlsx}).
      \item \textbf{Costificazione delle variazioni:} Possibilità per il Project Management di costificare le variazioni
            tecniche in modo separato dal contratto commerciale, così da valutare il margine di manovra sulle modifiche in corso d'opera.
\end{itemize}

\subsection{Gestione contrattuale, revisioni e ``As-Built''}
\begin{itemize}
      \item \textbf{Doppio contratto (Commerciale vs Tecnico):} Poiché il contratto commerciale firmato e sdoganato ha
            vincoli legali e non può essere alterato liberamente, il sistema deve gestire in parallelo un ``Contratto Tecnico''
            (copia speculare).
      \item \textbf{Allineamento Contratto vs Schema (P\&ID) e Feedback Loop:}
            \begin{itemize}
                  \item Sostituzione di note esterne e file Word con una sezione dedicata che registri le discrepanze tra
                        venduto commerciale ed eseguito tecnico (es. modifiche allo schema P\&ID, numero valvole, componenti extra),
                        mantenendo allineati i dati tecnici ed economici.
                  \item Implementazione di un sistema di feedback \textit{``As-Built''} per registrare le
                        modifiche effettive apportate durante la realizzazione, ottimizzando i processi e la precisione dei futuri preventivi.
                  \item Possibilità di riproporre al cliente il contratto tecnico aggiornato affinché diventi
                        il documento ufficiale.
            \end{itemize}
\end{itemize}

\subsection{Workflow, pianificazione temporale e notifiche}
\begin{itemize}
      \item \textbf{Monitoraggio date e colli di bottiglia:} Registrazione e confronto automatico
            tra la \textit{Data emissione ordine}
            (conferma cliente) e la \textit{Data stampa/presa in carico} (ricezione interna), evidenziando eventuali ritardi.
      \item \textbf{Pianificazione a ritroso (Back-Scheduling):} Calcolo automatico delle milestone intermedie a partire dalla
            \textit{Data di evasione} (spedizione prevista) o dallo \textit{Start-up impianto}
            per garantire il rispetto delle tempistiche.
      \item \textbf{Gestione stati e Alerting automatico:} Distinzione chiara degli stati della commessa
            (es. da \textit{Pre-inserimento / Bozza commerciale} a \textit{Inserito in Impianti / Validato tecnicamente}).
            Al cambio di stato, invio di notifiche automatiche (email o alert in-app) ai reparti a valle (es. Elettrico, Logistica),
            sostituendo il passaggio fisico del fascicolo cartaceo e le email manuali.
\end{itemize}

\section{Requisiti su anagrafiche, tracciabilità e governance}
\begin{itemize}
      \item \textbf{Tripla entità anagrafica e storicizzazione:}
            \begin{itemize}
                  \item Distinzione obbligatoria in anagrafica tra tre soggetti: \textit{Intestatario} (Fiscale),
                        \textit{Utilizzatore} (Fisico) e \textit{Destinatario Merce} (Logistico).
                  \item Storicizzazione dei dati anagrafici e introduzione di \textit{date di obsolescenza} per articoli
                        e sedi: in caso di cambio sede o P.IVA del cliente, i documenti storici (es. DDT passati) devono conservare
                        i dati validi all'epoca dell'emissione.
            \end{itemize}
      \item \textbf{Tracciabilità totale (Audit Log) e responsabilità:}
            \begin{itemize}
                  \item Eliminazione delle annotazioni su carta a favore di una visibilità digitale in tempo reale per tutti i
                        reparti autorizzati (Commerciale, \ac{UT}, \ac{PM}, Logistica, Magazzino).
                  \item Log completo delle modifiche ai dati di commessa che registri \textit{Chi} ha effettuato la variazione,
                        \textit{Quando} (timestamp) e \textit{Cosa} è cambiato (valore precedente vs valore attuale).
                  \item Presenza di un campo ``Compilato da'' autocompilato con l'utente loggato, ma mantenuto editabile
                        per gestire eventuali deleghe o correzioni di assegnazione.
            \end{itemize}
      \item \textbf{Governance Top-Down e formazione:}
            \begin{itemize}
                  \item Impostazione rigida del database: divieto di popolamento libero da parte degli utenti e definizione
                        delle regole a monte tramite un team ristretto per preservare l'integrità dei dati.
                  \item Formazione specifica per il reparto commerciale sull'utilizzo del nuovo sistema e sull'importanza
                        della correttezza dei dati inseriti a monte.
            \end{itemize}
\end{itemize}